% META: {"id": "TCOMPL-ATT-ME-005", "chapitre": "TCOMPL-TEMPS-ATTENTE", "type_objet": "methode", "capacites_codes": ["C5"], "methodes": ["M5"], "mode_creation": "ex_nihilo", "sources_inspiration": [], "status": "needs_review"}
% BEGIN-VERIFY
% from sympy import symbols, Rational, exp, integrate, oo, simplify
% t = symbols('t', positive=True)
% lam = Rational(1,2)
% f = lam*exp(-lam*t)
% assert integrate(f, (t, 0, oo)) == 1
% P13 = integrate(f, (t, 1, 3))
% assert simplify(P13 - (exp(Rational(-1,2)) - exp(Rational(-3,2)))) == 0
% assert 0 < float(P13) < 1
% END-VERIFY
\begin{fichemethode}{M5}{Verifier une densite de probabilite et calculer des probabilites}

\quandUtiliser{%
  Utiliser cette methode lorsque l'enonce donne une fonction $f$ et demande de
  VERIFIER que c'est une densite de probabilite, puis de calculer des
  probabilites du type $P(a \leq T \leq b)$ ou $P(T > a)$.
}

\pasApas{%
  \begin{enumerate}
    \item Verifier la POSITIVITE : $f(t) \geq 0$ sur l'intervalle ou elle est
      definie (justifier, ex. exponentielle $> 0$).
    \item Verifier la MASSE TOTALE : $\displaystyle\int f = 1$ sur le domaine
      (integrale sur $[0\,;\,+\infty[$ : calculer $\int_0^A f$ puis faire
      tendre $A$ vers $+\infty$).
    \item Conclure : $f$ positive et d'integrale 1 est une densite.
    \item Probabilites : $P(a \leq T \leq b) = \displaystyle\int_a^b f(t)\,dt$
      (les bornes larges ou strictes ne changent rien en continu).
    \item Pour $P(T > a)$ : utiliser le complementaire $1 - \int_0^a f$, ou la
      forme directe $e^{-\lambda a}$ pour une exponentielle.
  \end{enumerate}
}

\exempleRedige{%
  Verifier que $f(t) = \dfrac{1}{2}\,e^{-t/2}$ ($t \geq 0$) est une densite,
  puis calculer $P(1 \leq T \leq 3)$.

  \medskip
  \commentaireMarge{Positivite + integrale egale a 1.}%
  Pour tout $t \geq 0$, $f(t) > 0$. De plus
  $\displaystyle\int_0^A \dfrac{1}{2}e^{-t/2}\,dt
  = \left[-e^{-t/2}\right]_0^A = 1 - e^{-A/2} \xrightarrow[A \to +\infty]{} 1$.
  Donc $f$ est une densite de probabilite.

  \medskip
  \commentaireMarge{Integrale entre les deux bornes.}%
  $P(1 \leq T \leq 3) = \displaystyle\int_1^3 \dfrac{1}{2}e^{-t/2}\,dt
  = \left[-e^{-t/2}\right]_1^3 = e^{-1/2} - e^{-3/2} \approx 0{,}383$.
}

\pieges{%
  \begin{itemize}
    \item \textbf{Piege 1.} Oublier l'une des DEUX conditions (positivite ET
      integrale 1) : les deux sont necessaires.
    \item \textbf{Piege 2.} Erreur de signe dans la primitive :
      une primitive de $\lambda e^{-\lambda t}$ est $-e^{-\lambda t}$
      (controler en derivant).
    \item \textbf{Piege 3.} Traiter l'integrale impropre sans limite : ecrire
      $\int_0^A$ puis $A \to +\infty$, pas directement $\infty$ dans le
      crochet.
  \end{itemize}
}

\verifier{%
  Deriver la primitive pour retrouver $f$, controler que chaque probabilite est
  dans $[0\,;\,1]$, et verifier
  $P(T \leq a) + P(T > a) = 1$ sur un exemple.
}

\sEntrainer{\refExos{M5}}

\end{fichemethode}
